\section{Historical NV/NV Investigation}
\label{app:nvnv}

This appendix preserves the NV/NV work that led to the retained design. It
is intentionally separated from the principal result because those kernels
use a different P scale format, approximation policy, and sometimes a
different denominator. Their timings must not be mixed with the final NV/MX
headroom.

\subsection{Initial stabilized path}

The first numerically robust full-FP4 route followed the conventional order:
\begin{equation}
  z \rightarrow m \rightarrow e^{z-m} \rightarrow a_B
  \rightarrow s_{\mathrm{E4M3}} \rightarrow \mathrm{E2M1}.
\end{equation}
It computed an exact online row maximum, produced floating probabilities,
reduced an N32 block maximum, encoded an E4M3 scale, divided by that scale,
and packed E2M1. This path remained finite on downstream inputs but measured
about 0.1884\,ms at S4096/H24. The sequence was numerically conservative and
serial: scale selection could not start until floating probabilities existed.

For a global NV encode factor $G$, payload and scale can be expressed as
\begin{align}
  S_B(G) &= Q_{\mathrm{E4M3}}\!\left(G\frac{a_B}{6}\right), \\
  q_j(G) &= \Qe\!\left(\frac{G p_j}{S_B(G)}\right), \\
  \widehat p_j(G) &= \frac{S_B(G)}{G}q_j(G).
\end{align}
If $G$ multiplies both numerator and denominator consistently and no
rounding boundary changes, it cancels. It is not a free accuracy knob once
scale rounding, underflow, and saturation are considered. Sweeping $G$
without changing the shiftless estimator produced little benefit; adding
full stabilization to make the sweep meaningful restored the expensive
path.

\subsection{Code-directed polynomial development}

The key useful idea from NV/NV was to target E2M1 decisions directly. A
cubic fit to the base-two locations of E2M1 thresholds was
\begin{equation}
  \begin{aligned}
    p(x)={}&0.07839806x^3+0.28625049x^2 \\
           &+0.63145205x+0.99202336.
  \end{aligned}
\end{equation}
After folding the score and block scale into its coefficients, two values
could be evaluated with packed Horner operations:
\begin{align}
  r_1 &= \operatorname{FMA2}(x,a,b), \\
  r_2 &= \operatorname{FMA2}(r_1,x,c), \\
  y &= \operatorname{FMA2}(r_2,x,d).
\end{align}
This was faster than a full exponential and more accurate than early linear
fits. It also established that the fitting target should be
$\Qe(f(x))$, not $f(x)$ itself.

An intermediate SFU/ALU hybrid issued four native \code{EX2} pairs first and
filled their latency with twelve cubic or affine pairs. The native samples
also supplied an approximate denominator. This improved over pure-SFU and
pure-ALU variants on that schedule, but it coupled numerator and denominator
sampling to one distribution.

The later affine endpoint
\begin{equation}
  \ell(x)=\max(0,1.62330034x+0.92083546)
\end{equation}
reduced non-native work to one packed FMA per pair. That mechanism survives
in the retained NV/MX path, with policy-specific coefficients and an exact
represented denominator.

\subsection{Sampled denominator}

The NV/NV throughput path sampled eight of 32 scalar probability values in
each quarter. If $\mathcal J_q$ is the rotated sample set, it estimated
\begin{equation}
  \widetilde L_B = A_B\,4
  \sum_{j\in\mathcal J_q}2^{x_j}.
\end{equation}
The factor four compensates for sampling one quarter of the values. This
saved a complete denominator pass but could miss a concentrated maximum or
misestimate a nonuniform tail. Two- and three-word denominator reductions
were faster diagnostics but failed downstream: in a 20-image ViT gate they
gave respectively 0\% and 75\% top-1 agreement with BF16. The retained
NV/MX path instead sums all represented E2M1 payload words.

\subsection{Policy ladder and distribution dependence}

The NV/NV investigation exposed \code{fast}, \code{universal}, and
\code{hao-l2} style policies. \code{Fast} used the most aggressive
shiftless approximation. \code{Universal} restored scale guards and a more
reliable row reference. \code{Hao-l2} retained full stabilization. Their
rough latency hierarchy at S4096/H24 was approximately 0.10, 0.11--0.13,
and 0.188\,ms, respectively, depending on the exact generation of the
kernel.

The important result was not one latency number. Fast NV/NV could have good
cosine on Gaussian scores and still generate non-finite or highly distorted
model outputs when its row-reference assumptions changed. Stabilization
fixed those cases but consumed the desired speedup. Fixed-scale sweeps on
the shiftless estimator did not reproduce the accuracy of the fully
stabilized E4M3 control.

\subsection{Why NV/NV was not retained}

NV/NV contributed three ideas that remain: code-directed approximation,
early partial P publication, and the need to evaluate speed and downstream
accuracy together. It was superseded for four reasons:
\begin{enumerate}
  \item an E4M3 P scale without a robust global factor can round
  low-amplitude N32 blocks to zero; applying the factor correctly fixes the
  format-level underflow but restores work on the exposed P path;
  \item independent block-16 scale work does not align as naturally with the
  N32 software producer;
  \item the robust stabilized path was too slow, while the fastest shiftless
  path was distribution-sensitive;
  \item the HAO-derived 512-column layout and MXFP4 scale handoff provided a
  cleaner way to retain small P blocks without adding another score bank.
\end{enumerate}

The fixed-schedule format table remains in the main paper because it is a useful
control: NV/NV can outperform NV/MX on a specific synthetic metric. The
format-range and downstream experiments explain why that isolated result is
not sufficient to select the production route.

\subsection{B300 D64 sweep and range failure}
\label{app:nvnv-d64}

The D64 specialization provides a wider test of this distinction. Its B300
NV/MX route combines fused TMEM load/reduction with eight native-EX2 pairs and
eight affine pairs per score quarter. Every NV/MX output in the 24-shape
matrix is finite. The raw shiftless NV/NV control is often locally more
accurate because its E4M3 scale can fit a benign block more closely, but it is
finite on only \BThreeDsixtyfourNVNVFiniteCases\ of the 24 tested cases.

\begin{table}[htbp]
\centering
\caption{Selected B300 D64 format diagnostics. NV/MX is the retained finite
route. NV/NV timings are shown only to explain the rejected control and must
not be treated as production performance when the status is non-finite.}
\label{tab:b300-d64-diagnostic}
\scriptsize
\begin{adjustbox}{max width=\textwidth}
\begin{tabular}{rrrrrrrrrl}
\toprule
$H$ & $S$ & \multicolumn{4}{c}{NV/MX} & \multicolumn{3}{c}{raw NV/NV} & Status \\
\cmidrule(lr){3-6}\cmidrule(lr){7-9}
 & & ms & TFLOP/s & Cosine & Rel.-$L_2$ & ms & Cosine & Rel.-$L_2$ & \\
\midrule
12 & 4096 & 0.066528 & 775 & 0.955818 & 0.295719 & 0.070464 & 0.962281 & 0.278304 & finite \\
24 & 4096 & 0.093248 & 1105 & 0.955551 & 0.296653 & 0.099296 & 0.962063 & 0.279040 & finite \\
24 & 32768 & 4.601808 & 1434 & 0.956089 & 0.295010 & 4.842496 & 0.962065 & 0.279155 & finite \\
32 & 2048 & 0.039872 & 862 & 0.954553 & 0.299511 & 0.039968 & 0.961772 & 0.280428 & finite \\
32 & 32768 & 6.115360 & 1438 & 0.957303 & 0.291325 & 6.465520 & -- & -- & non-finite \\
64 & 1024 & 0.025568 & 672 & 0.952035 & 0.306387 & 0.025632 & 0.960558 & 0.286304 & finite \\
64 & 4096 & 0.207840 & 1323 & 0.955838 & 0.295956 & 0.218176 & 0.962128 & 0.278683 & finite \\
64 & 32768 & 12.223360 & 1439 & 0.957248 & 0.291501 & 12.919808 & 0.962800 & 0.276050 & finite \\
\bottomrule
%
\end{tabular}
\end{adjustbox}
\end{table}
\FloatBarrier

Across S4096-and-longer rows, retained NV/MX is
\BThreeDsixtyfourMinSpeedup--\BThreeDsixtyfourMaxSpeedup$\times$ faster than
the matched B300 BF16 kernel. Density two wins or ties density one on
\BThreeDsixtyfourDensityWins\ of 24 shapes. H32/S2048 and H64/S1024 need
136-CTA and 128-CTA grid caps, respectively, because their short logical job
counts interact poorly with the full 148-worker persistent grid. Other shapes
retain 148 CTAs. Every promoted build uses 128 registers, one barrier, and no
local-memory spills.

The decisive failure occurs at H32/S32768. Seed 20260802 produces exactly 64
non-finite raw NV/NV outputs while BF16 and NV/MX remain finite; another seed
passes. Saturating the E4M3 scale encoder repairs the failing distribution at
\BThreeDsixtyfourBoundedNVNVTime\,ms and
\BThreeDsixtyfourBoundedNVNVSpeedup$\times$ BF16, with 0.962921 cosine and
0.275742 relative-$L_2$. This bounded mode costs roughly 2--3\%. MXFP4's E8M0
scale is not finer than E4M3, but its exponent range represents tiny
probability blocks without a separately materialized row shift. Raw NV/NV can
therefore win a local error metric when it remains in range; NV/MX is the
retained low-latency route because it remains finite across the matrix.

\subsection{Measured downstream failure mechanism}
\label{app:nvnv-failure}

The expanded provider matrix isolates the failure. Native HAO and the TK
fixed-schedule control receive the same NVFP4-quantized Q, K, and V tensors.
HAO first subtracts a row maximum. The TK control instead forms a shiftless
block scale from
\begin{equation}
  s_B^{\mathrm{shiftless}}
  = \frac{1}{6}\max_{j\in B}\exp(z_j),
\end{equation}
then encodes it in E4M3. Model scores make this quantity exceed the finite
E4M3 maximum of 448 in every tested workload. This control has no
satfinite guard, so its out-of-range scale conversion contaminates the
dependent P and denominator path and produces non-finite context rows.

With row-max stabilization,
\begin{equation}
  s_B^{\mathrm{stable}}
  = \frac{1}{6}\max_{j\in B}\exp(z_j-m),
  \qquad m=\max_j z_j,
\end{equation}
so $0\le s_B^{\mathrm{stable}}\le 1/6$. Table~\ref{tab:nvnv-failure}
reports the first-sample diagnostic. The non-finite count is accumulated
over all attention layers evaluated for that sample. Small stabilized blocks
can still underflow, but they contain negligible probability mass relative
to the row anchor; HAO remains finite on every full fixed-input evaluation.

\begin{table}[htbp]
\centering
\caption{Shiftless TK NV/NV failure on model activations. Overflow is the
fraction of N32 P scales above E4M3's maximum before encoding.}
\label{tab:nvnv-failure}
\scriptsize
\begin{adjustbox}{max width=\textwidth}
\begin{tabular}{lrrrrrr}
\toprule
Task & Failed sample & Non-finite rows & Shiftless overflow (\%)
& Shiftless max & Stable overflow (\%) & Stable max \\
\midrule
ViT S256 & 1 & 2,560 & 8.146 & 1.828e+37 & 0.000 & 0.166667 \\
ViT S1024 & 1 & 127,950 & 8.137 & 5.332e+37 & 0.000 & 0.166667 \\
ViT S4096 & 1 & 539,442 & 7.511 & 5.669e+37 & 0.000 & 0.166667 \\
BERT MLM S256 & 1 & 167 & 1.721 & 5.279e+04 & 0.000 & 0.166667 \\
BERT MLM S512 & 1 & 49,613 & 1.063 & 5.655e+04 & 0.000 & 0.166667 \\
BERT SST-2 S256 & 1 & 21,717 & 1.465 & 4.493e+03 & 0.000 & 0.166667 \\
\bottomrule
%
\end{tabular}
\end{adjustbox}
\end{table}
